\section{Estratificación generacional, soportes quark y realización neutrínica}
\label{sec:generaciones-quarks-neutrinos-rev2}
\index{generaciones fermiónicas!estratificación HMT--MD}
\index{quark!soportes celulares y registro}
\index{PMNS!transporte de bases}

La ley general actúa sobre multisecciones y no sobre una lista de masas. La
aparición de tres generaciones se expresa, por ello, como compatibilidad de
tres estructuras: la estratificación del atlas, los registros de sabor y el
espacio tridimensional de mezcla. En el sector leptónico cargado, las
multisecciones
\begin{equation}
 G_0^{\ell},\qquad G_2^{\ell},\qquad G_4^{\ell}
 \label{eq:tres-multisecciones-leptonicas-rev2}
\end{equation}
tienen rangos \(1,8,16\). El centro \(G_0^\ell\) contiene la sección
electrónica de rango uno; \(G_2^\ell\) y \(G_4^\ell\) conservan, en cambio,
espectros operatoriales no triviales. La organización generacional completa
es
\begin{center}
\small
\begin{tabularx}{0.96\textwidth}{@{}c c c c X@{}}
\toprule
\textbf{Generación} & \textbf{Leptón} & \textbf{Par quark} &
\textbf{Sabor neutro} & \textbf{Soporte estructural}\\
\midrule
1 & \(e\) & \((u,d)\) & \(\nu_e\) & centro \(G_0^\ell\), registro de
primera generación y primer vector de la base de interacción\\
2 & \(\mu\) & \((c,s)\) & \(\nu_\mu\) & multisección \(G_2^\ell\),
registro de segunda generación y segundo vector de la base de interacción\\
3 & \(\tau\) & \((t,b)\) & \(\nu_\tau\) & multisección \(G_4^\ell\),
registro de tercera generación y tercer vector de la base de interacción\\
\bottomrule
\end{tabularx}
\end{center}
La conjugación partícula--antipartícula invierte las coordenadas de carga y
orientación de la representación, pero conserva esta estratificación; por
consiguiente, no duplica el número de generaciones.

La dimensión tres dispone además de un portador de incidencia propio. La
representación orientada dodecafásica contiene
\begin{equation}
 V_{12}\cong V_1\oplus V_3\oplus V_{3'}\oplus V_5,
 \qquad \operatorname{rank}P_3=3,
 \label{eq:descomposicion-rango-tres-generaciones-rev2}
\end{equation}
y el entrelazador dodecafásico--Witt transporta el sector positivo según
\begin{equation}
 P_G\xleftrightarrow{\;W_{GK}\;}P_3
 \xleftrightarrow{\;U_W\;}Q_3,
 \qquad
 \operatorname{rank}P_G=\operatorname{rank}P_3
 =\operatorname{rank}Q_3=3.
 \label{eq:transporte-rango-tres-generaciones-rev2}
\end{equation}
Este portador tridimensional permite comparar bases generacionales. Las
rutas, las firmas y los operadores de masa proporcionan después el contenido
espectral que CKM y PMNS transportan.

\subsection{Posición interna y registros de los \(6\) sabores de quark}

La posición de un sabor quark es el dato compuesto
\begin{equation}
 \mathfrak q_X=
 (\text{rama de carga},\text{generación},(\mathbb Z/3\mathbb Z)_{\rm color},
  \text{multisección},\text{registro de ruta}).
 \label{eq:posicion-quark-compuesta-rev2}
\end{equation}
Las ramas de carga \(-1/3\) poseen los soportes
\begin{align}
 D_1&=\{(4,5),(6,5)\},\nonumber\\
 D_2&=\{(4,3),(4,7),(6,3),(6,7)\},\nonumber\\
 D_3&=\{(2,3),(2,5),(2,7),(8,3),(8,5),(8,7)\},\nonumber\\
 D_4&=\{(2,1),(2,9),(4,1),(4,9),(6,1),(6,9),(8,1),(8,9)\},
 \label{eq:soportes-quark-down-rev2}
\end{align}
mientras las ramas de carga \(+2/3\) ocupan
\begin{align}
 U_1&=\{(4,4),(4,6),(6,4),(6,6)\},\nonumber\\
 U_3&=\{(2,2),(2,4),(2,6),(2,8),(4,2),(4,8),
       (6,2),(6,8),(8,2),(8,4),(8,6),(8,8)\}.
 \label{eq:soportes-quark-up-rev2}
\end{align}
Cada soporte se eleva a cuatro orientaciones por celda. El cociente de color
organiza \(36=12_{\rm R}+12_{\rm G}+12_{\rm B}\) biografías orientadas y conserva una misma
masa por sabor antes de la realización de color.

El transductor nominal
\begin{equation}
 \mathcal N_{\rm ruta}(\Gamma_f)
 =((N,E,S,O),h_{369},r_5)
 \in\mathbb Z_{\ge0}^{4}\times\mathbb Z_{\ge0}^{2}
 \label{eq:transductor-nominal-quark-rev2}
\end{equation}
publica las coordenadas siguientes antes de toda realización escalar:
\begin{center}
\small
\begin{tabularx}{0.96\textwidth}{@{}c c c c c X@{}}
\toprule
\textbf{Sabor} & \textbf{Generación} & \textbf{Rama} &
\textbf{\((N,E,S,O)\)} & \textbf{\((h_{369},r_5)\)} &
\textbf{Eje dominante}\\
\midrule
\(u\)&1&\(+2/3\)&\((34,31,31,12)\)&\((51,13)\)&N--S\\
\(d\)&1&\(-1/3\)&\((27,35,18,28)\)&\((47,17)\)&E--O\\
\(c\)&2&\(+2/3\)&\((20,38,24,26)\)&\((49,16)\)&E--O\\
\(s\)&2&\(-1/3\)&\((50,32,12,14)\)&\((57,11)\)&N--S\\
\(t\)&3&\(+2/3\)&\((46,27,16,19)\)&\((57,9)\)&N--S\\
\(b\)&3&\(-1/3\)&\((12,63,9,24)\)&\((47,7)\)&E--O\\
\bottomrule
\end{tabularx}
\end{center}
Las firmas nominales asociadas son
\begin{center}
\small
\begin{tabular}{@{}c c c@{}}
\toprule
\textbf{Sabor} & \textbf{\((n_A,s_C)\)} & \textbf{\((W_+,W_-)\)}\\
\midrule
\(u\)&\((11,7)\)&\((73,59)\)\\
\(d\)&\((17,8)\)&\((110,94)\)\\
\(c\)&\((61,8)\)&\((374,358)\)\\
\(s\)&\((41,-3)\)&\((243,249)\)\\
\(t\)&\((100,-1)\)&\((599,601)\)\\
\(b\)&\((71,-5)\)&\((421,431)\)\\
\bottomrule
\end{tabular}
\end{center}
Estas coordenadas distinguen sabores que comparten rama o generación y cierran
el morfismo de sabor sin consultar ninguna masa. Para una sección física (P),
definimos
\begin{equation}
 \operatorname{sabor}(P)
 =\bigl(\sigma_{\rm ud}(P),g(P),\chi_{\rm fina}(P)\bigr),
 \qquad
 \chi_{\rm fina}(P)
 =(\text{giros},\text{switches},\text{eje},N,E,S,O,h_{369},r_5),
 \label{eq:morfismo-sabor-ruta-cerrado-rev2}
\end{equation}
y, para quarks, se conserva además la órbita ternaria
\([\gamma_P]_{\rm col}=\{\gamma_P^R,\gamma_P^G,\gamma_P^B\}\). La realización
física es el mapa ya determinado por los datos precedentes,
\begin{equation}
 \mathcal S_{\rm phys}(P,\mathfrak M):
 \bigl(\sigma_{\rm ud},g,\chi_{\rm fina},[\gamma_P]_{\rm col}\bigr)
 \longmapsto \gamma^{\rm surv}_{P,\mathfrak M}
 \longmapsto
 \bigl(\sigma(\gamma^{\rm surv}_{P,\mathfrak M}),
       B_F^D,B_F^\Omega\bigr),
 \label{eq:realizacion-fisica-sabor-ruta-cerrada-rev2}
\end{equation}
donde \(\mathfrak M\) es el acoplamiento compatible y
\(\gamma^{\rm surv}_{P,\mathfrak M}\) la ruta que persiste bajo él. Para
leptones cargados la órbita de color se reduce al objeto terminal. El carácter
de masa se evalúa entonces sobre esa misma genealogía:
\begin{equation}
 \boxed{
 \frac{m_{\mathfrak M}(P)}{m_e^{\rm HMT}}
 =\chi_{\rm mass}\!\left(
   \sigma\!\left(\gamma^{\rm surv}_{P,\mathfrak M}\right)\right).}
 \label{eq:masa-como-caracter-persistencia-rev2}
\end{equation}
Por tanto, sabor, generación y color no esperan un selector futuro: publican la
clase de ruta y su autovector antes de abrir la comparación externa. La masa no
elige retrospectivamente la ruta; sólo lee multiplicativamente la persistencia
que el acoplamiento ha estabilizado.

\subsection{Fermionicidad, autovalores de masa y transporte PMNS}

La clausura exterior del funtor de intercambio construye las relaciones CAR
y fija la fermionicidad. En el sector neutro, la etiqueta
\(\nu_e,\nu_\mu,\nu_\tau\) designa una base de interacción; los valores de
masa son autovalores de un operador neutro mezclado. Los marcos internos
\(F_\ell,F_\nu:\mathbb C^3\to\operatorname{im}P_3\) comparan esas bases, y
la matriz
\begin{equation}
 U_{\rm HMT}=F_\ell^*\Phi_LF_\nu
 \label{eq:transporte-pmns-generaciones-rev2}
\end{equation}
es un transporte unitario. PMNS cambia de base entre estados de interacción
y autovectores de masa; la escala neutra pertenece al espectro del operador,
no al cambio de base.

\subsection{Extractor neutro Z0 anterior a la proyección}
\label{sec:extractor-neutro-z0-torsion-rev2}

La construcción recuperada del sector neutro actúa sobre el registro
dodecafásico enriquecido, antes de colapsar hoja, torsión y bit
bidireccional. Sean \(E_m\), \(m\in\mathbb Z/12\mathbb Z\), los doce
eventos de nueve ticks del registro CT--108. En cada tick \(t\) se conservan
\[
 \theta(t)=\mathbf1_{\{d(t)=9\}},\qquad
 \varepsilon(t)\in\{-1,+1\},\qquad
 \kappa(t)=\mathbf1_{\{9\ {\rm en\ corona}\}},
 \qquad \omega(t)=(-1)^{\kappa(t)}.
\]
El primer factor detecta la torsión nonádica, el segundo conserva la hoja
two--ways y el tercero cambia el signo en el giro de \(180^\circ\) de la
corona. El entero torsional de cada evento y su diferencia antipodal son
\begin{equation}
 T_\nu(m)=\sum_{t\in E_m}
   \omega(t)\varepsilon(t)\theta(t),\qquad
 e_m^\nu=T_\nu(m)-T_\nu(m+6),
 \label{eq:extractor-z0-torsion-antipodal-rev2}
\end{equation}
y el trit neutro queda fijado, sin parámetros continuos, por
\begin{equation}
 \tau_m^\nu=\operatorname{sgn}(e_m^\nu)
 \in\{-1,0,+1\}.
 \label{eq:trit-neutro-torsion-antipodal-rev2}
\end{equation}

Para las cuatro órbitas de salto cuatro
\[
 I_a=\{a,a+4,a+8\}\pmod{12},\qquad a=0,1,2,3,
\]
el coeficiente de evento que entra en la firma y el observable de mayoría
de signos son, respectivamente,
\begin{equation}
 b_{120}^\nu
 =\sum_{a=0}^{3}\operatorname{sgn}
   \left(\sum_{m\in I_a}e_m^\nu\right),
 \qquad
 \nu_{120}^{\nu,\mathrm{sgn}}
 =\sum_{a=0}^{3}\operatorname{sgn}
   \left(\sum_{m\in I_a}\tau_m^\nu\right).
 \label{eq:canales-120-neutros-distintos-rev2}
\end{equation}
No se identifican estos dos enteros: el primero es la coordenada
\(b_{120}=\nu_{120}^{\rm ev}\) del carácter, mientras el segundo es un
diagnóstico de mayoría. La memoria cuadrática es
\begin{equation}
 b_\zeta^\nu=\nu_{270}^\nu
 =\sum_{m=0}^{11}\tau_m^\nu\tau_{m+3}^\nu.
 \label{eq:canal-270-neutro-rev2}
\end{equation}
Por tanto, para cada historia neutra superviviente
\(\gamma_{\nu_i}\), el registro publica antes de toda comparación la firma
\[
 v_\nu(\gamma_{\nu_i})
 =\bigl(0,0,k_{\Delta_4}^{(i)},b_{120}^{\nu_i},
        b_\zeta^{\nu_i}\bigr)
\]
y el exponente relativo
\begin{equation}
 \mathcal A_\nu^{(i)}
 =k_{\Delta_4}^{(i)}\Delta_4
  +b_{120}^{\nu_i}s_{120}
  +b_\zeta^{\nu_i}\zeta_{270},
 \qquad
 \zeta_{270}=135\Delta_4^2.
 \label{eq:caracter-relativo-neutro-z0-rev2}
\end{equation}
Aquí \(\zeta_{270}\) no se sustituye por el momento de torre
\(s_{270}=q_-^{270}-q_+^{270}\).

La conjugación global two--ways
\(\tau^\nu\mapsto-\tau^\nu\) invierte
\(b_{120}^\nu\) y conserva \(b_\zeta^\nu\). En consecuencia genera dos
ramas de espectro relativo,
\begin{equation}
 \mathcal A_{\nu,\pm}^{(i)}
 =k_{\Delta_4}^{(i)}\Delta_4
  \pm b_{120}^{\nu_i}s_{120}
  +b_\zeta^{\nu_i}\zeta_{270},
 \qquad
 \text{NO}\longleftrightarrow\text{IO},
 \label{eq:bifurcacion-neutra-no-io-rev2}
\end{equation}
que se distinguen por el signo posterior de
\(\Delta m_{31}^2\), no por una selección con datos PDG. Esta construcción
cierra el extractor discreto Z0 y la bifurcación relativa. La escala absoluta
común \(m_0\) y el núcleo físico PMNS de rango completo permanecen objetos
separados: no se los introduce dentro del selector ni se los declara
deducidos por la sola paridad two--ways.

La realización con una única fase de rango uno produce
\[
 (\theta_{12},\theta_{23},\theta_{13})
 =(19.6807^\circ,56.4970^\circ,29.6224^\circ)
\]
y no reproduce la fenomenología de mezcla. Este control selecciona como dominio adecuado un
núcleo de registro de rango completo,
\begin{equation}
 G_{ab}=\frac1{\sqrt{|L_a||N_b|}}
 \sum_{c\in L_a}\sum_{d\in N_b}
 K_{\rm registro}(c,d;\xi_c,\xi_d),
 \qquad U_{\rm int}=\operatorname{polar}(G),
 \label{eq:kernel-registro-rango-completo-rev2}
\end{equation}
cuya fuente técnica se desarrolla en el capítulo de transporte CKM y
PMNS. Así quedan separados el operador de masa, sus autovalores y el mapa de
lectura por sabor.
